\section{Likelihoods, intervals and profiling}\label{app:C.14}

A search that sees no events with no expected background excludes, at \SI{90}{\percent} confidence, every signal expectation above $N_{90} = 2.303$ (\cref{eq:1.5-N90}). Real searches add backgrounds, several bins and uncertain inputs, and \cref{sec:8.8,sec:21.4,sec:21.5} handle them with likelihoods: the Poisson and Gaussian likelihoods, the likelihood ratio and the intervals it defines, nuisance parameters and their profiling, and, in \cref{sec:17.6}, Monte Carlo propagation of input widths. The meaning of a confidence level, a statement about the procedure in repeated experiments, is that of \cref{sec:1.5}.

\subsection{Poisson and Gaussian likelihoods}

The \emph{likelihood} $\mathsf{L}$ is the probability of the observed data, read as a function of the parameters. A counting experiment that observes $N_{\mathrm{obs}}$ events, with expected signal $N_{\mathrm{sig}}$ and background $N_{\mathrm{bkg}}$, has by \cref{eq:1.5-poisson}
\begin{equation}
  \mathsf{L}(N_{\mathrm{sig}}) = \frac{(N_{\mathrm{sig}} + N_{\mathrm{bkg}})^{N_{\mathrm{obs}}}}{N_{\mathrm{obs}}!}\,\ee^{-(N_{\mathrm{sig}} + N_{\mathrm{bkg}})} .
  \label{eq:C.14-poisson}
\end{equation}
Independent bins multiply, so a spectrum in bins $j$ has the product of factors of this form, one per bin, with expectations $N_{\mathrm{sig},j} + N_{\mathrm{bkg},j}$~\cite[Eq.~(6)]{cowan2011asymptotic}. Only ratios of likelihoods enter what follows, and factors independent of the parameters, such as $1/N_{\mathrm{obs}}!$, drop out.

When a parameter is measured with a Gaussian error, for example an estimate $\hat N$ of $N_{\mathrm{sig}}$ with standard deviation $\sigma_{\hat N}$ from many events, the likelihood is Gaussian,
\begin{equation}
  \mathsf{L}(N_{\mathrm{sig}}) \propto \exp\Bigl[-\frac{(\hat N - N_{\mathrm{sig}})^2}{2\sigma_{\hat N}^2}\Bigr] .
  \label{eq:C.14-gauss}
\end{equation}
The Poisson likelihood takes this form at large counts. Expanding $\ln\mathsf{L}$ of \cref{eq:C.14-poisson} to second order about its maximum, at $N_{\mathrm{sig}} + N_{\mathrm{bkg}} = N_{\mathrm{obs}}$, gives $\ln\mathsf{L} = \mathrm{const} - (N_{\mathrm{sig}} + N_{\mathrm{bkg}} - N_{\mathrm{obs}})^2/(2N_{\mathrm{obs}})$, a Gaussian of width $N_{\mathrm{obs}}^{1/2}$. The cubic term of the expansion is smaller by a factor of order $N_{\mathrm{obs}}^{-1/2}$ at one standard deviation from the maximum, so the Gaussian form needs tens of events, not a few.

\subsection{The likelihood ratio and intervals}

The \emph{likelihood ratio} compares a hypothesized value with the best one,
\begin{equation}
  \mathsf{t}(N_{\mathrm{sig}}) = -2\ln\frac{\mathsf{L}(N_{\mathrm{sig}})}{\mathsf{L}(\hat N_{\mathrm{sig}})} \ge 0 ,
  \label{eq:C.14-ratio}
\end{equation}
where $\hat N_{\mathrm{sig}}$ maximizes $\mathsf{L}$~\cite[Eq.~(8)]{cowan2011asymptotic}. Large $\mathsf{t}$ means poor agreement. For the Gaussian likelihood \cref{eq:C.14-gauss}, $\hat N_{\mathrm{sig}} = \hat N$ and $\mathsf{t} = (\hat N - N_{\mathrm{sig}})^2/\sigma_{\hat N}^2$ exactly. If the data are repeated, $\hat N$ scatters as a Gaussian about the true value, and $\mathsf{t}$ at the true value is the square of a standard normal variable, a $\chi^2$ variable with one degree of freedom. The set of $N_{\mathrm{sig}}$ with $\mathsf{t}(N_{\mathrm{sig}}) \le z^2$, the interval $\hat N \pm z\sigma_{\hat N}$, then covers the true value with probability $2\Phi(z) - 1$, with $\Phi$ the standard normal distribution function. For a likelihood that is not Gaussian, Wilks' theorem states that $\mathsf{t}$ at the true value approaches the same $\chi^2$ distribution as the sample grows~\cite{wilks1938large}, \cite[Sect.~3.1]{cowan2011asymptotic}, and the same rule applies. The thresholds are
\begin{equation}
  \mathsf{t} \le 1.00,\ 2.71,\ 3.84 \qquad \text{for } \SI{68.27}{\percent},\ \SI{90}{\percent},\ \SI{95}{\percent}\ \text{(two-sided)},
  \label{eq:C.14-thresholds}
\end{equation}
the squares of $z = 1$, 1.645 and 1.960. A $1\sigma$ interval is the \SI{68.27}{\percent} interval, written ``\SI{68}{\percent}'' throughout the book.

An upper limit is one-sided. Data with $\hat N_{\mathrm{sig}}$ above the tested value are not counted as evidence against it, so the statistic is set to zero there~\cite[Eq.~(14)]{cowan2011asymptotic}, and all of the rejected probability lies on one side. For a Gaussian likelihood the upper limit at confidence level $\mathrm{CL}$ is $\hat N + z_{\mathrm{CL}}\sigma_{\hat N}$ with $\Phi(z_{\mathrm{CL}}) = \mathrm{CL}$: $z = 1.282$ at \SI{90}{\percent} and $z = 1.645$ at \SI{95}{\percent}, which correspond to $\mathsf{t} = 1.64$ and 2.71. The same number, 2.71, thus marks a two-sided \SI{90}{\percent} interval and a one-sided \SI{95}{\percent} limit.

The asymptotic rule fails at small counts, the regime of dark-matter limits. With $N_{\mathrm{obs}} = 0$ and $N_{\mathrm{bkg}} = 0$, \cref{eq:C.14-poisson} gives $\mathsf{L} = \ee^{-N_{\mathrm{sig}}}$, maximal at $\hat N_{\mathrm{sig}} = 0$, and $\mathsf{t} = 2N_{\mathrm{sig}}$. The one-sided \SI{90}{\percent} threshold 1.64 would give $N_{\mathrm{sig}} < 0.82$, while the exact construction of \cref{sec:1.5} gives 2.303. Limits from a few events must therefore be built from the exact Poisson probabilities, by the classical construction of \cref{sec:1.5} or by its generalization that orders outcomes by the likelihood ratio and yields one- or two-sided intervals as the data demand~\cite{feldman1998unified}.

\subsection{Nuisance parameters and the profile likelihood}

The predicted counts depend on quantities that are not of interest but are not known exactly: a background normalization, an efficiency, or the size of a truncation in the theory (\cref{sec:8.8}). Each becomes a \emph{nuisance parameter} $\vartheta$ in the likelihood, usually with a constraint term that encodes what is known about it from a calibration or from a computed next order. The \emph{profile likelihood} replaces each nuisance parameter by its best value at fixed $N_{\mathrm{sig}}$,
\begin{equation}
  \mathsf{t}_{\mathrm{p}}(N_{\mathrm{sig}}) = -2\ln\frac{\mathsf{L}\bigl(N_{\mathrm{sig}}, \hat{\hat\vartheta}(N_{\mathrm{sig}})\bigr)}{\mathsf{L}(\hat N_{\mathrm{sig}}, \hat\vartheta)} ,
  \label{eq:C.14-profile}
\end{equation}
where $\hat{\hat\vartheta}(N_{\mathrm{sig}})$ maximizes $\mathsf{L}$ at fixed $N_{\mathrm{sig}}$ and the denominator is the overall maximum~\cite[Eq.~(7)]{cowan2011asymptotic}. Wilks' theorem holds for $\mathsf{t}_{\mathrm{p}}$ as well, with one degree of freedom per parameter of interest, so \cref{eq:C.14-thresholds} still applies.

A Gaussian example shows what profiling does. Let the estimate be shifted by an unknown offset, so that $\hat N$ is Gaussian with mean $N_{\mathrm{sig}} + \vartheta$ and width $\sigma_{\hat N}$, and let a calibration constrain the offset by a Gaussian of mean 0 and width $\sigma_\vartheta$. Then
\begin{equation*}
  -2\ln\mathsf{L}(N_{\mathrm{sig}}, \vartheta) = \frac{(\hat N - N_{\mathrm{sig}} - \vartheta)^2}{\sigma_{\hat N}^2} + \frac{\vartheta^2}{\sigma_\vartheta^2} + \mathrm{const} .
\end{equation*}
Write $r = \hat N - N_{\mathrm{sig}}$. Setting the derivative with respect to $\vartheta$ to zero, $-(r - \vartheta)/\sigma_{\hat N}^2 + \vartheta/\sigma_\vartheta^2 = 0$, gives $\hat{\hat\vartheta} = r\sigma_\vartheta^2/(\sigma_{\hat N}^2 + \sigma_\vartheta^2)$. Inserting it, the two terms become $r^2\sigma_{\hat N}^2/(\sigma_{\hat N}^2 + \sigma_\vartheta^2)^2$ and $r^2\sigma_\vartheta^2/(\sigma_{\hat N}^2 + \sigma_\vartheta^2)^2$, and their sum is $r^2/(\sigma_{\hat N}^2 + \sigma_\vartheta^2)$. The overall minimum is zero, at $N_{\mathrm{sig}} = \hat N$ and $\vartheta = 0$. Hence
\begin{equation}
  \mathsf{t}_{\mathrm{p}}(N_{\mathrm{sig}}) = \frac{(\hat N - N_{\mathrm{sig}})^2}{\sigma_{\hat N}^2 + \sigma_\vartheta^2} .
  \label{eq:C.14-quadrature}
\end{equation}
The profiled interval is that of a Gaussian of width $(\sigma_{\hat N}^2 + \sigma_\vartheta^2)^{1/2}$: statistical and systematic widths add in quadrature. A nuisance parameter that scales the normalization, $N_{\mathrm{sig}} \to N_{\mathrm{sig}}(1 + \vartheta)$, gives the same result for the relative widths to first order in $\vartheta$. Profiling broadens the likelihood, and the broadening is the information about $N_{\mathrm{sig}}$ lost to the nuisance parameter. A nuisance parameter that the data constrain better than the calibration does is pulled towards the value the data prefer, which a fixed-width envelope would miss.

\subsection{Monte Carlo propagation}

When a derived quantity depends on many inputs through a formula that is not linear, its distribution is obtained by sampling. Each input $x_j$ is given a distribution, for example lognormal, with $\ln x_j$ Gaussian of mean $\ln x_j^{(0)}$ and width $\sigma_j$, a width set by how the input was obtained. We draw $N_{\mathrm{MC}}$ sets of inputs, evaluate the output $y$ on each, and read off percentiles. The 16th and 84th percentiles bound the central \SI{68}{\percent} interval, because $\Phi(-1) = 0.159$; for a Gaussian output they are the mean minus and plus one standard deviation, and for any output they define an interval of the same coverage that need not be symmetric.

Comparisons follow from the same draws. The probability that material $A$ outperforms $B$ is the fraction of draws in which $y_A > y_B$. For lognormal outputs with log-widths $\sigma_A$, $\sigma_B$ and correlation $r$ between their errors, $\ln y_A - \ln y_B$ is Gaussian with mean $\Delta = \ln y_A^{(0)} - \ln y_B^{(0)}$ and variance $\sigma_A^2 + \sigma_B^2 - 2r\sigma_A\sigma_B$, and
\begin{equation}
  P(A \succ B) = \Phi\Bigl(\frac{\Delta}{(\sigma_A^2 + \sigma_B^2 - 2r\sigma_A\sigma_B)^{1/2}}\Bigr) .
  \label{eq:C.14-pairwise}
\end{equation}
With $\Delta = 0.3$, $\sigma_A = 0.2$, $\sigma_B = 0.25$ and $r = 0$, $P(A \succ B) = \Phi(0.94) = 0.83$. A positive correlation, as arises when both materials are computed with the same method and share its error, shrinks the variance of the difference and sharpens the comparison; a ratio of two such outputs is more certain than either. The sampling error of an estimated probability $p$ is $[p(1 - p)/N_{\mathrm{MC}}]^{1/2}$, at most 0.0035 for $N_{\mathrm{MC}} = 2\times10^4$. All such probabilities are conditional on the input distributions: they propagate the assigned widths faithfully, but they cannot be more reliable than the widths themselves.
